%% Stand 12. Juli 2026
\subsection{Kriging-Gleichungen}
\label{KapitelTheorieKriging}

Unter Kriging subsummiert man eine Klasse geostatistischer
Schätzverfahren, deren Ziel die möglichst genaue Vorhersage
einer räumlichen (oder als Erweiterung auch raum-zeitlichen) 
Zufallsvariable an nicht beobachteten Orten ist. 
Im Gegensatz zu anderen Interpolationsverfahren werden 
hierbei sowohl die räumliche Korrelation als auch die
Unsicherheit der Schätzung explizit berücksichtigt.

Die Herleitung der Kriging-Gleichungen basiert auf drei grundlegenden Prinzipien:
\begin{itemize}[topsep=2pt,parsep=1pt,itemsep=0.5pt]
\item[$\circ$] Die Vorhersage erfolgt als lineare Kombination der Beobachtungen.
\item[$\circ$] Die Schätzung soll \underline{erwartungstreu} 
sein. Das heißt für einen Schätzer $\theta$ der gesuchten 
Größe $T$
soll $\mathbb{E}\left( \theta\right) = T$ gelten. \\
Dies ist gleichbedeutend zu der Forderung einer 
\underline{unverzerrten} Schätzung,
d.\,h.\,mit $b_{\theta} = (\mathbb{E}\left( \theta \right) - T)$ soll gelten, dass 
$b_{\theta} = 0$ erfüllt ist
(vgl.\,z.\,B.\,\cite[S.\,235]{Henze2019}).
\item[$\circ$] Unter allen unverzerrten linearen Schätzern wird
diejenige gewählt, deren Fehlervarianz minimal ist.
\end{itemize}

Die verschiedenen Kriging-Verfahren unterscheiden sich im Wesentlichen durch 
die Annahmen über den Erwartungswert des zugrunde liegenden Zufallsfeldes. 
Im Folgenden werden zunächst die gemeinsame Grundidee und anschließend die 
Herleitung der drei Varianten Simple, Ordinary und Universal Kriging erläutert.

\subsubsection{Methodik des Kriging}
Es sei eine Zufallsvariable $\mathrm{Z}(s)$ mit 
$s\in \mathcal{S}$ gegeben. Gesucht ist eine Vorhersage 
am Ort $s_0$.

Die Kriging-Schätzung wird als gewichtete Summe der Beobachtungen angesetzt:
\begin{equation}
\hat{\mathrm{Z}}(s_0) =
\sum_{i=1}^{n}\lambda_i \mathrm{Z}(s_i),
\end{equation}
wobei $\lambda_i$ die noch unbekannten Gewichte darstellen.\\
Intuitiv sollen die Gewichte so gewählt werden, 
dass Beobachtungen mit hoher räumlicher Ähnlichkeit stärker berücksichtigt werden 
als weiter entfernte Beobachtungen mit geringerer Ähnlichkeit.
Gleichzeitig sollen Redundanzen vermieden werden: 
Liegen mehrere Messpunkte sehr dicht beieinander, liefern sie nur
begrenzte zusätzliche Information und erhalten daher geringere
 Einzelgewichte.\\
Die Bestimmung der Gewichte erfolgt über ein Optimierungsproblem. 
Hierzu wird die mittlere quadratische Vorhersageabweichung
\begin{equation}
\sigma_{\mathrm{opt}} =
\mathrm{Var}
\left(
\hat{\mathrm{Z}}(s_0)- \mathrm{Z}(s_0)\right)
\end{equation}
minimiert, wobei zusätzlich die Bedingung der Erwartungstreue 
erfüllt werden muss.

Die Varianz des Schätzfehlers hängt ausschließlich von der
Kovarianzstruktur beziehungsweise äquivalent von der Semivarianz 
des Zufallsfeldes ab. Dadurch fließen die räumlichen 
Abhängigkeiten unmittelbar in die Berechnung der Gewichte ein.

\subsubsection{Simple Kriging}
Beim \underline{Simple Kriging} wird angenommen, dass der Erwartungswert des Zufallsfeldes 
bekannt und im gesamten Untersuchungsgebiet konstant ist,

\begin{equation}
\mathbb{E}\left( \mathrm{Z}(s) \right)=m.
\end{equation}

Diese Annahme vereinfacht die Herleitung erheblich, da keine zusätzlichen Bedingungen 
an die Gewichte notwendig sind.\\
Die resultierende Schätzung kann geschrieben werden als:
\begin{equation}
\hat{\mathrm{Z}}(s_0) = m + \sum_{j=1}^{n} \lambda_j
\left( \mathrm{Z}(s_j)- m \right),
\end{equation}
wobei $n$ die Anzahl der Beobachtungen bezeichnet.\\ 
Somit werden nicht die absoluten Beobachtungswerte interpoliert, sondern deren 
Abweichungen vom bekannten Mittelwert $m$.

Die Gewichte werden so gewählt, dass die Varianz des Schätzfehlers minimal wird. 
Da der Mittelwert bekannt ist, genügt hierfür die Lösung eines linearen Gleichungssystems, 
welches ausschließlich die Kovarianzen zwischen den Beobachtungspunkten sowie 
die Kovarianzen zum Vorhersageort enthält.

Formal ergibt sich das Gleichungssystem
\begin{equation}
\mathbf{C}\cdot \lambda = c,
\end{equation}
wobei

\begin{itemize}[topsep=2pt,parsep=1pt,itemsep=0.5pt]
\item[$\circ$] $\mathbf{C}$ die Kovarianzmatrix der Beobachtungen,
\item[$\circ$] $c$ der Kovarianzvektor zwischen Beobachtungen 
und Vorhersageort sowie
\item[$\circ$] $\lambda$ der Vektor der unbekannten Gewichte
\end{itemize}
symbolisieren. Für eine eindeutige Lösbarkeit des
Gleichungssystems ist es erforderlich, dass die Kovarianzmatrix $\mathbf{C}$ positiv definit sein muss.\\

Simple Kriging besitzt damit die mathematisch einfachste Form aller Kriging-Verfahren. 
In praktischen Anwendungen ist die Voraussetzung eines bekannten globalen 
Mittelwertes allerdings selten erfüllt. Aus diesem Grund wird Simple Kriging 
vergleichsweise selten eingesetzt und dient hauptsächlich als 
theoretischer Ausgangspunkt.

\subsubsection{Ordinary Kriging}
\label{KapitelTheorie_OrdinaryKriging}
In den meisten Anwendungen ist der Mittelwert des Prozesses unbekannt. 
Stattdessen wird lediglich angenommen, dass dieser innerhalb einer zu
spezifizierenden Nachbarschaft $N\subset \mathcal{S}$ konstant ist, also
\begin{equation}
\mathbb{E}\left( \mathrm{Z}(s) \right) = m,\ s \in N,
\end{equation}
wobei $m$ unbekannt bleibt.\\
Die große Bedeutung der
Nachbarschaft wird z.\,B.\,in \cite[Kap.\,8.5]{Webster2007} 
diskutiert, wo auch empirische Regeln zu deren Festlegung 
aufgeführt sind. \\
Da der Mittelwert nicht bekannt ist, muss verhindert werden, dass die 
Schätzung systematisch verzerrt wird. Dies führt unmittelbar zur 
Nebenbedingung

\begin{equation}
\sum_{j=1}^{n}\lambda_j = 1.
\end{equation}

Diese Bedingung stellt sicher, dass sich der unbekannte 
Mittelwert bei der Bildung des Erwartungswertes herauskürzt 
und die Schätzung unverzerrt bleibt.

Die Optimierung erfolgt wiederum über die Minimierung der
Schätzfehlervarianz. Aufgrund der zusätzlichen Nebenbedingung 
 kann das Optimierungsproblem jedoch nicht mehr direkt gelöst 
werden. Stattdessen wird die Methode der Lagrange-Multiplikatoren
verwendet.\\
Hierzu wird die Zielfunktion $L$ um den Lagrange-Multiplikator 
$\xi$ erweitert:

\begin{equation}
L = \sigma_{\mathrm{opt}} + 2\xi \left( \sum_j\lambda_j - 1 \right).
\end{equation}

Durch Ableiten nach allen Gewichten und dem Lagrange-Multiplikator 
$\xi$ entsteht das Kriging-Gleichungssystem

\begin{equation}
\begin{bmatrix}
\mathbf{C} & \mathbf{1}\\
\mathbf{1}^T & 0
\end{bmatrix}
\begin{bmatrix}
\lambda\\
\xi
\end{bmatrix}
=
\begin{bmatrix}
c\\
1
\end{bmatrix}.
\end{equation}

Im Vergleich zum Simple Kriging wird somit lediglich eine zusätzliche Gleichung 
ergänzt, welche die Unverzerrtheit garantiert.

Ordinary Kriging stellt den am häufigsten verwendeten Kriging-Ansatz dar, da 
keine Kenntnis des globalen Mittelwertes erforderlich ist und dennoch 
optimale lineare Schätz\-un\-gen erzeugt werden.

\subsubsection{Universal Kriging}
Während Ordinary Kriging einen lokal konstanten Mittelwert voraussetzt, reicht 
diese Annahme in vielen Anwendungen nicht aus.\\ 
Universal Kriging zerlegt den Prozess daher in einen 
deterministischen Trend und einen räumlich korrelierten Restprozess 
$\epsilon(s)$,
also
\begin{equation}
\mathrm{Z}(s) = m(s) + \varepsilon(s).
\end{equation}

Dabei wird der Trend als Linearkombination bekannter 
Basisfunktionen beschrieben
\begin{equation}
m(s) = \sum_{k=1}^{p} \beta_k f_k(s),
\end{equation}
wobei
\begin{itemize}[topsep=2pt,parsep=1pt,itemsep=0.5pt]
\item[$\circ$] $f_k(s)$ bekannte Trendfunktionen  (z.\,B.\,Polynomfunktionen) und
\item[$\circ$] $\beta_k$ unbekannte Regressionsparameter
\end{itemize}
bezeichnen. Der Trend wird also durch ein multiples lineares 
Regressionsmodell beschrieben. 
Die Schätzbedingung lautet weiterhin

\begin{equation}
\label{Gl_Schätzbedinung_UniversalKriging}
\hat{\mathrm{Z}}(s_0) = \sum_{j=1}^n \lambda_j \mathrm{Z}(s_j),
\end{equation}

jedoch müssen nun zusätzlich die $p$ Trendfunktionen reproduziert 
werden. Daraus ergeben sich die Nebenbedingungen

\begin{equation}
\sum_j \lambda_j f_k(s_j) =
f_k(s_0), \qquad 1 \leq k \leq p.
\end{equation}

Im Vergleich zum Ordinary Kriging existiert somit nicht 
nur eine, sondern mehrere Nebenbedingungen.

Auch hier erfolgt die Herleitung über Lagrange-Multiplikatoren. Das resultierende 
lineare Gleichungssystem besitzt die Form

\begin{equation}
\begin{bmatrix}
\mathbf{C} & \mathbf{F}\\
\mathbf{F}^T & \mathbf{0}
\end{bmatrix}
\begin{bmatrix}
\lambda\\
\xi
\end{bmatrix}
=
\begin{bmatrix}
c\\
\mathbf{f}(s_0)
\end{bmatrix},\ \xi \in \mathbb{R}^{(p+1)},
\end{equation}

wobei $\mathbf{F}$ die Matrix der Trendfunktionen enthält.

Universal Kriging kann daher als Verallgemeinerung des Ordinary Kriging 
verstanden werden. Während Ordinary Kriging lediglich einen konstanten 
Mittelwert zulässt, erlaubt Universal Kriging beliebige lineare Trends.

\subsubsection{Zusammenfassender Überblick}
Alle drei Methoden fußen auf denselben 
Grundanahmen einer linearen, erwartungstreuen und
varianzminimalen Prädiktion mit zunehmend allgemeineren 
Annahmen an den Mittelwert.

\begin{itemize}[topsep=2pt,parsep=1pt,itemsep=0.5pt]
\item[$\circ$] Simple Kriging: 
\begin{itemize}[topsep=2pt,parsep=1pt,itemsep=0.5pt]
\item[$\triangleright $] $m$ bekannt und konstant; 
\item[$\triangleright $] keine Nebenbedingung; 
\item[$\triangleright $] unrestringierte Optimierung.
\end{itemize}
\item[$\circ$] Ordinary Kriging: 
\begin{itemize}[topsep=2pt,parsep=1pt,itemsep=0.5pt]
\item[$\triangleright $] $m$ unbekannt, aber regional konstant;
\item[$\triangleright $] eine Nebenbedingung ($\sum_j \lambda_j = 1$) und somit
\item[$\triangleright $] ein Lagrange-Multiplikator.
\end{itemize}
\item[$\circ$] Universal Kriging: 
\begin{itemize}[topsep=2pt,parsep=1pt,itemsep=0.5pt]
\item[$\triangleright $] $m(s)$ unbekannt und variabel (Trendfunktion); 
\item[$\triangleright $] $(p+1)$ Nebenbedingungen und somit 
\item[$\triangleright $] $(p+1)$ Lagrange-Multiplikatoren.
\end{itemize}
\end{itemize}

Diese Hierarchie erklärt zugleich, warum sich die drei Verfahren in ihrer 
mathematischen Struktur so ähnlich sind: Das jeweilige Kriging-System 
entsteht stets durch dieselbe Lagrange-Konstruktion, lediglich die Anzahl 
und Form der Nebenbedingungen wächst mit der Komplexität der an
den Mittelwert gestellten Annahmen und erhöht folglich die Dimension 
des zu lösenden Gleichungssystems.
%Anhand der Gl.\,(\ref{Gl_Schätzbedinung_UniversalKriging})
%erkennt man relativ leicht, warum Kriging zur Klasse der
%\textit{Interpolationsverfahren} gehört. Angenommen, es sei $s_0 = s_j$ für ein 
%$1 \leq j \leq n$.\\
%Dann folgt, dass
%\begin{equation}
%\hat{\mathrm{Z}}(s_j) = \sum_{k=1}^n \lambda_k \mathrm{Z}(s_k) 
%\ \mathrm{mit\ } \lambda_j =1\ \wedge\ \lambda_k = 0\,
%\forall\,k \neq j
%\end{equation}
%eine \textit{mögliche} Lösung darstellt, welche die
%Nebenbedingungen erfüllt. Da unter der Voraussetzung einer
%positiv definiten Kovarianz-Matrix die Lösung der 
%Kriging-Gleichungen eindeutig ist, folgt, dass 
%$\hat{\mathrm{Z}}(s_j)= \mathrm{Z}(s_j)$ gelten muss.
%Einschränkend sei bemerkt, dass im  Falle der Einbeziehung 
%eines Nugget-Effekts die obige Beweis-Skizze \underline{nicht} mehr
%gilt; in diesem Fall ist der Kriging-Prädiktor kein perfekter Interpolator
%mehr.
%\\
%Im raum-zeitlichen Kontext überträgt sich diese Systematik unmittelbar, 
%indem die Ko\-va\-ri\-anz-Funktion sowie gegebenenfalls die Trendfunktionen 
%$f_k$ als Funktionen des kombinierten Raum-Zeit-Index $\mathbf{x} = (s,t)$ 
%definiert werden.